\FloatBarrier
\section{System II: Edwards Aquifer, index well J-17}\label{sec:edwards}

\subsection{Data and specification}\label{sec:edw_data}

The predictand is the calendar-year mean of daily-high elevation at well J-17 (TWDB 6837203, EAA AY-68-37-203), the index well of the San Antonio Pool of the Edwards (Balcones Fault Zone) Aquifer, for 1934--2023 (Table~\ref{tab:edwspec}, Fig.~\ref{fig:edwseries}). Three hydrological objects are read off this one number and should be kept apart: the head is observed at an access point; it indexes a store of water, which is the resource; and the store is replenished and drawn by fluxes, which are the flow. The one-pool water balance closes this loop approximately, in the increment-structure sense used throughout: head change equals weighted fluxes plus a linear drain. It is not a closed mass balance. Annual means use all available daily highs; no year falls below the 240-observation floor (minimum $n=242$ in 1939); missing days are not interpolated; the published pre-continuous composite is used as issued; and the 2023 values carry provisional status. Recharge $R$ is USGS total San Antonio-area recharge ($10^3$ acre-ft\,yr$^{-1}$), estimated by the stream-loss method of \citet{Puente1978} and constructed independently of J-17 \citep{UmphresChoi2025}; pumpage $P$ is well discharge from Edwards Aquifer Authority Table 1 \citep{EAA2024}. Both are constructed fluxes, not observations of J-17, and total-area values mix the San Antonio and Uvalde pools. Spring discharge is stored and not used as a driver. Climate predictors, used only by the climate-informed variants, are September--November Ni\~no~3.4 (HadISST) and calendar-year precipitation in Texas climate divisions 06 and 07 (NCEI nClimDiv). Two thresholds frame the series without entering any forecast: the Edwards Aquifer Authority's Stage~I line (10-day mean J-17 below 660\,ft, a 2007 institutional rule) and the physical level near 618\,ft at which Comal Springs approaches cessation (1956 daily minimum 612.51\,ft). The annual mean is below 660\,ft in 31 of 90 years; the daily minimum is below 618\,ft in one (1956).

\begin{table}[!htb]
\centering\small
\caption{The J-17 specification. D = data, E = empirical construct, M = model, N = normative threshold. The predictand is a measured well series, not a storage reconstruction, a numerical-model inversion, the J-27 Uvalde index or spring discharge.}
\label{tab:edwspec}
\begin{tabularx}{\textwidth}{@{}>{\raggedright\arraybackslash}p{0.24\textwidth}Lc@{}}
\toprule
Field & Contents & Type \\
\midrule
System & Edwards Aquifer, San Antonio Pool, indexed by J-17 & D \\
Predictand $H_t$ & calendar-year mean of daily-high J-17 elevation (ft AMSL), 1934--2023 (90 years) & D \\
Drivers & recharge $R_t$ (USGS, stream-loss method); pumpage $P_t$ (EAA Table 1) & D/E \\
Physical threshold & head at which Comal Springs does not cease ($\approx$618\,ft) & N/E \\
Institutional threshold & EAA Stage I, 10-day mean $<660$\,ft (in force after 2007); the annual-mean indicator is a coarse proxy & N \\
Disturbance & recharge pulses; unmodelled karst; Uvalde--San Antonio exchange & M \\
Service series & USGS 08168710 Comal Springs annual-mean discharge (cfs); scored after retention is frozen & D \\
Fixed windows & drought-of-record drawdown (train 1934--1950, test 1951--1956); drought-of-record recovery (1934--1956 / 1957--1961); pre-permit wet interval (1980--1990 / 1991--1995); critical-period era (1997--2014 / 2015--2023) & --- \\
Rolling origins & minimum 15 training years; $n=75$ ($h=1$), $n=71$ ($h=5$); every model on identical origin sets & --- \\
\bottomrule
\end{tabularx}
\end{table}

\begin{figure}[!htb]
\centering
\includegraphics[width=0.92\textwidth]{edw_fig1_series}
\caption{J-17 annual mean and daily-high range, 1934--2023. 1956 mean 623.15\,ft (daily minimum 612.51\,ft); 1992 mean 691.96\,ft (daily maximum 703.31\,ft); 2023 mean 635.68\,ft. The annual mean is below the 660-ft Stage I line in 31 of 90 years.}
\label{fig:edwseries}
\end{figure}

\subsection{The water-balance ladder}\label{sec:edw_ladder}

\begin{definition}[One-pool water-balance map]\label{def:pool}
With head clipped to $[610,710]$\,ft, the one-pool map is
\begin{equation}\label{eq:pool}
H_{t+1} \;=\; \bigl[\,H_t + \alpha + \beta\tilde R_{t+1} + \gamma\tilde P_{t+1} + \delta H_t\,\bigr]_{\mathrm{clip}},
\end{equation}
where $\tilde R$ and $\tilde P$ are the flux values a forecast substitutes for the unknown year-$(t+1)$ fluxes. Forecasts are issued at the end of year $t$ from $(H,R,P)$ through $t$; $(\alpha,\beta,\gamma,\delta)$ are refitted by least squares on the training transitions at every origin.
\end{definition}

The clip stands in for a physical storage floor; it never binds on the observed record (annual means 623--692\,ft) and binds only on persisted-drought forecast paths, where it does real work (SI S17). The ladder (Table~\ref{tab:edwladder}) differs from the cod ladder in one structural respect that Section~\ref{sec:synthesis} turns on: the causal rungs must \emph{substitute} something for the target-year fluxes. M2 persists last year's $(R_t,P_t)$; M2m substitutes the training-mean fluxes; M3 adds an AR(1) residual on $\Delta H$ to M2; M4 is a one-year information-delay variant kept as a symmetry control with the cod ladder (J-17 is a telemetered gauge, so the module prices a theoretical lag rather than a constraint of this system). The climate-informed variants replace persisted $R$ by a one-step forecast of $R$ from information known at $t$: an AR(1) on $R$, lagged Ni\~no~3.4, lagged climate-division precipitation, or all three; for $h>1$ the one-step recharge forecast is held constant. The oracle receives the realised $(R_{t+h},P_{t+h})$ and is excluded from retention, as is a precipitation oracle. M1 is the univariate AR(1), $H_{t+1}=a+\varphi H_t$, scored as the autonomous rung.

\begin{table}[!htb]
\centering\small
\caption{The J-17 ladder. ``Comparator'' is the next-simpler causal model under (E2); M1 has none and is tested against persistence only.}
\label{tab:edwladder}
\begin{tabularx}{\textwidth}{@{}>{\raggedright\arraybackslash}p{0.13\textwidth}>{\raggedright\arraybackslash}p{0.16\textwidth}>{\raggedright\arraybackslash}p{0.2\textwidth}L>{\raggedright\arraybackslash}p{0.14\textwidth}@{}}
\toprule
ID & Class & Fluxes substituted at $t+h$ & Role & Comparator \\
\midrule
persist & baseline & --- & $\hat H_{t+h}=H_t$ & --- \\
mean & baseline & --- & training-window mean & --- \\
M1 & autonomous AR(1) & --- & $H_{t+1}=a+\varphi H_t$ & (persistence only) \\
M2 & causal stock-flow & last $(R_t,P_t)$ persisted & candidate for retention & M1 \\
M2m & climatological stock-flow & training-mean $(\bar R,\bar P)$ & affine AR(1) by class & M1 \\
M3 & residual & as M2 & AR(1) residual on $\Delta H$ & M2 \\
M4 & delay & as M2 & information lag; symmetry control & M3 \\
M2\_Rar, M2\_Renso, M2\_Rprecip, M2\_combo & climate-informed & one-step forecast of $R$ (AR(1) on $R$; Ni\~no~3.4; precipitation; all three) & climate rung & M1 (declared) \\
M2\_oracle & diagnostic & realised $(R,P)$ & nowcast ceiling; excluded & --- \\
\bottomrule
\end{tabularx}
\end{table}

\begin{remark}[Class reduction of M2m]\label{rem:class}
Under constant fluxes $\tilde R_{t+1}=\bar R$, $\tilde P_{t+1}=\bar P$, the map \eqref{eq:pool} reduces to $H_{t+1}=(1+\delta)H_t+\mathrm{const}$ and therefore shares M1's forecast function class, though not its estimator: M2m pins its intercept from the in-sample mean fluxes, an additional identifying use of the recharge and pumpage records, while the persistence coefficient $(1+\delta)$ is not mean-flux dependent. The decision to decline M2m on these grounds is a protocol choice recorded in the frozen documents, not a consequence of (E1)--(E2).
\end{remark}

\FloatBarrier
\subsection{The verdict under Protocol E, and under Protocol C}\label{sec:edw_verdict}

Table~\ref{tab:edwrolling} gives the rolling-origin summary and Fig.~\ref{fig:edwrmse} the per-model RMSE at both horizons. At $h=1$ persistence scores 13.23\,ft and the training mean 16.17\,ft. M1 scores 12.84\,ft and satisfies (E1); as the simplest autonomous model it has no next-simpler causal comparator, so (E2) is vacuous and M1 is retained by the point rule. M2 (14.70), M3 (14.46) and M4 (14.30\,ft) each fail (E1): the causal balance that persists last year's fluxes is worse than persisting the head. M2m (12.28\,ft) satisfies (E1) and, against M1, (E2); it is listed and then declined by the class clause. The oracle, at 7.55\,ft, is 43\% below persistence and is excluded. Under Protocol E, therefore, one output-only module is retained and no stock-flow, residual or delay module is. Under Protocol C, restated on the same scores, nothing is retained: M1's 2.96\% margin lies inside the 5\% band, and at $h=5$ M1 (21.25\,ft) does not beat persistence (21.11\,ft) at all. The two protocols disagree about M1 and agree about everything else, and both readings are reported because both are true of the scores.

\begin{table}[!htb]
\centering\small
\caption{J-17: rolling-origin summary (ft; $n=75$ at $h=1$, $n=71$ at $h=5$, identical origin sets for every row) and the retention decision under Protocol E. Subscripts give the horizon; $\Delta$ columns are $h=1$ RMSE differences, negative favouring the module. The training mean is the best non-oracle forecaster at $h=5$.}
\label{tab:edwrolling}
\footnotesize\setlength{\tabcolsep}{3.5pt}
\begin{tabularx}{\textwidth}{@{}>{\raggedright\arraybackslash}p{0.19\textwidth}rrrrrL@{}}
\toprule
Model & RMSE$_1$ & MAE$_1$ & RMSE$_5$ & $\Delta$ persist & $\Delta$ M1 & Protocol E decision \\
\midrule
persist & 13.23 & 10.73 & 21.11 & --- & --- & benchmark \\
mean & 16.17 & 13.17 & \textbf{16.80} & $+2.94$ & --- & benchmark \\
M1 AR(1) & 12.84 & 10.72 & 21.25 & $-0.39$ & --- & retained (point rule; margin within noise) \\
M2 persisted fluxes & 14.70 & 11.45 & 33.49 & $+1.47$ & $+1.86$ & rejected: fails (E1) \\
M2m climatological fluxes & 12.28 & 10.22 & 17.44 & $-0.95$ & $-0.56$ & passes (E1), (E2); declined by class clause \\
M3 residual & 14.46 & 11.12 & 33.46 & $+1.23$ & $+1.62$ & rejected \\
M4 delay & 14.30 & 11.17 & 33.39 & $+1.07$ & $+1.46$ & rejected \\
M2\_oracle & 7.55 & 5.79 & 10.86 & $-5.68$ & --- & excluded (uses future $R,P$) \\
\bottomrule
\end{tabularx}
\end{table}

\begin{figure}[!htb]
\centering
\includegraphics[width=0.8\textwidth]{edw_fig3_rmse}
\caption{J-17: rolling-origin RMSE by model at $h=1$ and $h=5$. The oracle is a diagnostic ceiling and is excluded from retention. At five years the training mean is the best non-oracle forecaster and M2--M4 are 60\% worse than persistence.}
\label{fig:edwrmse}
\end{figure}

The retention of M1 is a one-year, point-RMSE statement and should be read as such. The margin is 0.39\,ft on $n=75$; MAE is a tie (10.72 against 10.73\,ft); at five years the ordering reverses; and at the decision scale of annual drought-stage declarations the difference is operationally nil. The record of the class clause is as important as the retention: on the rule as written, the statement ``no stock-flow module is retained'' is an artefact of declining M2m, since the water-balance-identified affine map with climatological fluxes is the best one-step forecaster tested. What the ladder actually shows is that \emph{persisting last year's recharge} is the failure, not the balance as such; Section~\ref{sec:edw_mech} and Proposition~\ref{prop:increment} make that statement exact.

\subsection{Uncertainty on the margins}\label{sec:edw_uncertainty}

Table~\ref{tab:edwdm} attaches Diebold--Mariano statistics and moving-block bootstrap intervals to the load-bearing margins, computed from the archived per-origin forecast files; the layer is post-freeze and changes no verdict. Three readings follow. The retention margin itself (M1 $-$ persist, $-0.39$\,ft) is within noise, as are the M2m-over-M1 estimator margin and the M2 causal loss (the last borderline, its interval touching zero); the only one-year margin that separates from zero is M2m over persistence ($-0.95$\,ft, $p=0.003$). The five-year climatology win ($-4.30$\,ft) is robust: its bootstrap interval excludes zero, and the $h=5$ Diebold--Mariano $p$ of 0.10 reflects the known under-sizing of that statistic under overlapping origins. And the climate-rung margins are within noise in both directions. An independent second implementation of the layer (different seed, block length $h$, 20\,000 replicates, different HAC weighting and variance convention; SI S18) reproduces every reading: M1's retention margin has $z=-0.86$ and block interval $[-1.22,+0.56]$\,ft; the M2m edge over persistence has $p=0.001$; the five-year climatology win has $p=0.035$.

\begin{table}[!htb]
\centering\small
\caption{J-17: uncertainty on the load-bearing margins (post-freeze layer; Diebold--Mariano with Newey--West HAC at lag $h-1$; paired moving-block bootstrap, block length 8, 10\,000 replicates, seeded). Rolling origin, $h=1$ unless noted.}
\label{tab:edwdm}
\scriptsize\setlength{\tabcolsep}{3pt}
\begin{tabular}{@{}lrrrcl@{}}
\toprule
Comparison & margin (ft) & DM & $p$ & 95\% block CI (ft) & separates \\
\midrule
M1 $-$ persist & $-0.39$ & $-0.85$ & 0.40 & $[-1.51,\,+0.71]$ & no \\
M2m $-$ persist & $-0.95$ & $-3.07$ & 0.003 & $[-1.45,\,-0.68]$ & \textbf{yes} \\
M2m $-$ M1 & $-0.56$ & $-1.61$ & 0.11 & $[-1.53,\,+0.13]$ & no \\
M2 $-$ persist & $+1.47$ & $+1.27$ & 0.21 & $[-0.02,\,+3.40]$ & no (borderline) \\
mean $-$ persist ($h=5$) & $-4.30$ & $-1.65$ & 0.10 & $[-9.72,\,-2.09]$ & \textbf{yes} \\
M1 $-$ persist ($h=5$) & $+0.15$ & $+0.06$ & 0.95 & $[-5.57,\,+4.94]$ & no \\
M2\_combo $-$ M1 (climate) & $-0.13$ & $-0.21$ & 0.83 & $[-1.37,\,+1.31]$ & no \\
M2\_combo $-$ M2m (nested) & $+0.43$ & $+0.74$ & 0.46 & $[-0.56,\,+1.57]$ & no \\
\bottomrule
\end{tabular}
\end{table}

\FloatBarrier
\subsection{The mechanism: a near-white driver, a nowcasting map, and what climate indices add}\label{sec:edw_mech}

The one-pool increment tracks same-year recharge: $r(\Delta H_t,R_t)=0.74$ on the 1934--2023 record, and the oracle's RMSE is 43\% below persistence at $h=1$ and 49\% below at $h=5$. The same increment is almost unrelated to last year's recharge ($r(\Delta H_t,R_{t-1})=0.06$), because annual recharge has almost no memory: $r(R_t,R_{t-1})=0.17$ against 0.87 for pumpage. A module that persists $R_t$ therefore persists a quantity that does not persist, and it loses to persisting the head itself; given the year's realised fluxes the same map is a nowcast of the year's head at 7.55\,ft. That is the diagnostic that separates a forecasting failure from a structural one: the balance accounts for the contemporaneous increment, and it has nothing to say about next year's increment because next year's recharge is not in any information set at the origin.

The fixed windows make the point visibly (SI S17). The 1957--1961 recovery from the drought of record is a recharge pulse ($R_{1956}=43.7$, $R_{1957}=1142.6\times10^3$ acre-ft): persistence stays at the 1956 floor (RMSE 43.6\,ft), M2 persists drought recharge and falls further, into the clip (55.3\,ft), and the oracle tracks the rise (12.3\,ft). The 1992 peak and the 2015--2023 window repeat the split. On the 1951--1956 drawdown M2 ``beats'' the oracle (18.1 against 19.7\,ft), but for a reason that undermines it: the map trained on 1934--1950 has the wrong sign on pumpage ($\gamma=+0.021$), because pumping rose as the drought deepened and least squares aliases the human response into the state transition.

The climate-informed variants ask whether the signals observable at the annual origin (September--November Ni\~no~3.4, lagged climate-division precipitation, an AR(1) on $R$) supply a recharge forecast that closes the gap (Table~\ref{tab:edwclimate}). They do not. On the recharge target itself the climate modules have modest skill relative to climatology (528--545 against $556\times10^3$ acre-ft; persisted recharge 702), and the autoregression on near-white recharge has none. On the head, the three climate modules beat persistence and M1 by 0.02--0.13\,ft, margins within noise (Table~\ref{tab:edwdm}); at $h=5$ all three are 3--6\,ft worse than persistence, a design consequence of holding the one-step recharge forecast constant over the horizon; and against the nested climatological-flux baseline M2m they lose outright (12.71--12.82 against 12.28\,ft). None is retained; the fixed-window record, including the one window on which two climate modules edge past M1 by about a foot, is in SI S22. Same-year precipitation, which is not available at the origin, would reduce head RMSE to 10.56\,ft (precipitation oracle), still short of the full flux oracle because the linear rain map misses 1957-scale extremes. The gap between persistence and the oracle is next year's recharge, and nothing known on 31 December recovers it.

\begin{table}[!htb]
\centering\small
\caption{J-17: climate-informed recharge, rolling origin on the same origin sets as Table~\ref{tab:edwrolling}. Head RMSE in ft; recharge RMSE in $10^3$ acre-ft. Margins are against M1, the declared comparator of the climate rung; M2m is the nested climatological-flux baseline against which all climate modules lose.}
\label{tab:edwclimate}
\begin{tabular}{@{}lrrrr@{}}
\toprule
Model & $H$, $h{=}1$ & margin vs.\ M1 & $H$, $h{=}5$ & $R$, $h{=}1$ \\
\midrule
persist $H$ / persist $R$ & 13.23 & $+0.39$ & \textbf{21.11} & 702 \\
M1 & 12.84 & --- & 21.25 & --- \\
M2\_Rar & 13.25 & $+0.41$ & 25.38 & 561 \\
M2\_Renso & 12.82 & $-0.02$ & 24.42 & 528 \\
M2\_Rprecip & 12.80 & $-0.04$ & 25.38 & 545 \\
M2\_combo & 12.71 & $-0.13$ & 26.88 & 538 \\
M2m (declined nested baseline) & 12.28 & $-0.56$ & 17.44 & --- \\
recharge climatology & --- & --- & --- & 556 \\
precipitation oracle (year $t{+}h$) & 10.56 & --- & 16.91 & \textbf{354} \\
\bottomrule
\end{tabular}
\end{table}

\subsection{Post-freeze layers: the service series and the pumpage counterfactuals}\label{sec:edw_post}

Two layers were scored after retention was frozen and change no verdict; both are reported in full in the SI. The Comal Springs discharge series is a near-linear channel of the same state ($r=0.986$ with J-17): the same ladder scored on it over the same origins reproduces the head ranking (persist 71.9, M1 69.0, M2m 68.7, M2 74.8\,cfs; oracle 45.3), as it must for a linear transform, and the one place the channel is not redundant with the head is the drought tail, where the fitted rating curve predicts about 97\,cfs at the 1956 mean against an observed 32 and never predicts cessation where cessation occurred (SI S20). The pumpage counterfactuals pass alternative pumpage paths through the pre-permit map from the 1990 head with the actual recharge sequence: the full policy spread, from a 20\% cut to freezing at the 1990 peak, spans 630.9--646.8\,ft at 2023, which is of the same order as the map's own RMSE against the observed record (7.2--14.2\,ft), so pumpage is a secondary lever on this map and no scenario repairs the ladder's primary failure (SI S21).

\subsection{Reading the aquifer verdict}\label{sec:edwreading}

Three statements survive the uncertainty layer. First, the causal balance that persists last year's recharge is not a one-year forecaster of J-17 head, and the reason is in the driver's spectrum rather than in the map; with the year's fluxes the map nowcasts. Second, the AR(1)'s retention under Protocol E is a weak-autocorrelation phenomenon, not evidence that an output-only model is generally preferable. An estimated AR(1) nests persistence at $\varphi=1$; as the head series approaches that limit the two forecasts converge and the AR(1)'s advantage shrinks toward the noise in estimating $\varphi$. On this record the head is persistent but far from a unit root (lag-one autocorrelation 0.63; fitted $\hat\varphi=0.66$), which leaves the AR(1) a real but slim margin, and the directional prediction is that on a more strongly autocorrelated head record the AR(1) should beat persistence by \emph{less}. An exploratory port of the ladder to the separately regulated Uvalde Pool (annual head autocorrelation 0.84), where the AR(1) loses to persistence by 0.13\,ft, is consistent with that direction and is offered as a direction only; no Uvalde result enters any verdict here. The fitted $\hat\varphi=0.66$ also admits a physical reading: to first order spring discharge obeys Darcy proportionality to head above the spring level, whose autonomous solution is an affine autoregression $H_{t+1}\approx(1-k)H_t+kH_s$ with drainage coefficient $\hat k\approx0.34$\,yr$^{-1}$; the module is output-only in the protocol's sense but the shape it fits carries the aquifer's free-drainage momentum. Third, at five years the training mean beats persistence robustly (16.80 against 21.11\,ft, interval excluding zero), the specification's most directly transferable result for multi-year planning in a rapidly recharged, institutionally bounded aquifer; a non-stationary fossil aquifer under sustained depletion would not sustain a mean-reverting baseline, and Section~\ref{sec:synthesis} gives the criterion that decides which case a basin is in.
